\subsection{Landslide susceptibility modeling}
\label{sec:related:landslide}

Statistical and machine-learning approaches to landslide susceptibility have been the subject of extensive review \citep{reichenbach2018review, merghadi2020ml}. The field has progressively moved from bivariate and weights-of-evidence approaches to logistic regression, random forests, gradient-boosted trees (XGBoost, CatBoost), and more recently to deep neural networks. \citet{brenning2005spatial} established spatial cross-validation as a methodological best practice for landslide modeling, demonstrating that random hold-out validation systematically overestimates generalization error due to spatial autocorrelation between nearby training and test pixels. Object-oriented approaches \citep{stumpf2011oo} and patch-based convolutional models have refined the spatial feature representation, but the fundamental requirement of hundreds to thousands of labeled events per study area remains a binding constraint. Recent benchmarks comparing classical and ensemble ML algorithms \citep{merghadi2020ml} report $F_1$ scores in the 0.75--0.90 range when $N>200$, but performance degrades sharply for small inventories. Most of the field operates within a single basin or contiguous region; cross-basin transfer is studied less frequently and is typically framed as out-of-sample testing rather than as a problem requiring methodological adaptation.

\subsection{Transfer learning in remote sensing}
\label{sec:related:transfer}

Transfer learning by fine-tuning a model pre-trained on data-rich source domains is the de facto strategy for label-efficient remote sensing \citep{tuia2016domain}. In landslide applications, \citet{ghorbanzadeh2022landslide} demonstrated that ImageNet-pretrained convolutional networks fine-tuned on a few hundred local samples outperform train-from-scratch baselines, confirming that representations learned on natural images contain transferable low-level features. However, fine-tuning a deep network with $K<50$ target labels typically results in catastrophic over-fitting unless aggressive regularization or early stopping is applied. Furthermore, the ``concat-source'' variant---pre-training on a union of source basins, which is the baseline we adopt---implicitly assumes that all source domains are equally relevant to the target, a hypothesis that is rarely tested explicitly and that fails when the source pool spans heterogeneous environmental conditions. Recent work in RSE has begun to formalize cross-region transferability through spatiotemporal deep-learning frameworks that integrate environmental context for agronomic mapping \citep{luo2026transferable}. Within the ISPRS Journal of Photogrammetry and Remote Sensing community, \citet{wei2024universaladapter} propose a universal adapter for transferable landslide segmentation, and \citet{ng2025crossresolution} address cross-resolution landslide mapping with deep learning. Adversarial domain-adaptation methods continue to evolve in the same venue \citep{zhang2025awda}, and semi-supervised event detection from multi-temporal optical imagery offers a complementary path for sparse-inventory settings \citep{deijns2024semisupervised}. However, landslide-specific cross-basin meta-learning with rigorous spatial validation remains an open question.

\subsection{Domain adaptation}
\label{sec:related:da}

Domain adaptation (DA) seeks to learn classifiers that generalize across distribution shifts by leveraging both labeled source data and unlabeled target data. The Domain-Adversarial Neural Network (DANN) of \citet{ganin2015dann} is the canonical method: a feature extractor is trained simultaneously to (i) discriminate the class label and (ii) confuse a domain classifier through a Gradient Reversal Layer, encouraging the learned representation to be domain-invariant. Variants based on Maximum Mean Discrepancy \citep{long2015dan} and Wasserstein distance have followed. In remote sensing, DA has been applied to multi-source land-cover classification, image segmentation, and scene recognition, with mixed success \citep{tuia2016domain}: performance gains are typically substantial when source domains are abundant and labeled target data are completely absent, but degrade when only a handful of source domains are available or when small numbers of labeled target samples are accessible. The latter regime is precisely our setting, and our results below confirm that DANN is not a competitive baseline when source-domain count is small.

\subsection{Meta-learning}
\label{sec:related:meta}

Meta-learning frames the data-scarcity problem as one of \emph{learning to learn}: the model is trained on a distribution of related tasks so that adaptation to a new task from few examples requires only a small parameter update. Model-Agnostic Meta-Learning (MAML; \citealt{finn2017maml}) is the canonical second-order method, computing gradients through inner adaptation steps, but it is computationally expensive and unstable in practice. First-order approximations---FOMAML, which drops the second-order term, and Reptile \citep{nichol2018reptile}, which replaces query-loss-driven updates with averaged inner-step displacements---are more practical. Theoretical work \citep{raghu2020rapid} suggests that MAML's success can be largely attributed to feature reuse rather than rapid weight adaptation, which we revisit in our adaptation-curve analysis (Section~\ref{sec:results}).

In remote sensing, meta-learning has been applied to few-shot land-cover classification \citep{russwurm2020metalearning}, achieving 70--90\% accuracy with $K=5$ samples per class on EuroSAT and Sen12MS benchmarks, and to satellite image classification more broadly \citep{tseng2022timm}. \citet{ma2025metalsm} recently demonstrated temporal meta-learning for dynamic landslide susceptibility mapping at a single site (Hong Kong), decomposing yearly subtasks and achieving fast adaptation in 5 samples and 5 inner steps. Our work complements this in two ways: (i) we focus on \emph{spatial} cross-basin transferability rather than temporal adaptation; (ii) we evaluate on 14 basins across a 1500~km climate gradient, providing the first cross-basin benchmark of first-order meta-learning algorithms (Reptile, FOMAML) for landslide susceptibility under rigorous spatial cross-validation.
